% =========================
% Cas pratique
% =========================

\subsection{Cas pratique}

%--------------------------
\begin{frame}{Cas pratique}

\begin{block}{}
\pctart{22}(1):
délai pour l'entrée en phase nationale DO/EP: 30 mois à compter de la date de priorité.
\end{block}

En l'espèce,
\begin{itemize}
  \item Priorité : 22/11/2022
  \item Dépôt PCT : 20/11/2023 (priorité valable Art. 4C CUP)
  \item Entrée en phase régionale EP envisagée : 31/03/2025
\end{itemize}

Donc, le délai d'entrée en phase nationale DO/EP
expire fin mai 2025.\newline

\textbf{Conclusion: Le 31/03/2025, le déposant
est dans le délai d'entrée en phase nationale
pour DO/EP.}

\end{frame}

%--------------------------
\begin{frame}{Cas pratique}

\begin{block}{}
\pctart{22}(1) + \pctrule{49} ; règles 159 et s. CBE:
\begin{itemize}
  \item requête d’entrée en phase ;
  \item paiement de la taxe nationale ;
  \item remise d’une traduction si nécessaire.
\end{itemize}
\end{block}

Demande en français (langue officielle OEB) donc pas de traduction à fournir.\newline

\textbf{Conclusion: il reste à déposer la requête d'entrée en phase et paiement de la taxe nationale.}

\end{frame}

%--------------------------
\begin{frame}{Cas pratique}

\begin{block}{}
\pctart{27}(2) + \pctrule{51bis}.1.a.vi): L’irrégularité n’empêche pas l’entrée en phase, mais l’OEB peut exiger une
confirmation / régularisation de la signature du déposant non signataire.
\end{block}

En l'espèce, deux déposants désignés pour EP (CA-P et FR-O),
mais FR-O n'a pas signé.\newline

\textbf{Conclusion : demander la signature de FR-O pour la régularisation
à l’entrée en phase.}

\end{frame}

%--------------------------
\begin{frame}{Cas pratique}

\begin{block}{}
\pctart{17}.3.b):  les parties de la demande internationale qui n'ont par conséquent pas fait l'objet d'une recherche 
sont considérées comme retirées pour ce qui concerne les effets dans cet état, 
à moins qu'une taxe particulière ne soit payée par le déposant à l'office national dudit état.
\end{block}

En l'espèce, l’ISA/CA a soulevé un défaut d’unité (3 inventions). CA-P a payé une taxe additionnelle
uniquement pour la 2\textsuperscript{e} invention, sans réserve la recherche effectuée
pour inventions 1 et 2 seulement. Toutefois, l’appréciation ISA n’engage pas l’OEB sur l’unité.\newline

\textbf{Conclusion: courrier de demande d'instruction client
pour anticiper si défaut d'unité retenu:}
\begin{itemize}
  \item option A: (i) argumenter, puis le cas échéant (ii) limiter et (iii) déposer une divisionnaire EP.
  \item option B: abandonner invention 3 sans argumenter.
\end{itemize}

\end{frame}

%--------------------------
\begin{frame}{Cas pratique}

\begin{block}{}
\pctart{27}(1) (et en phase EP : art. \textbf{123}(2) CBE / EPC).
\end{block}

En l'espèce, les dessins ne figurent ni dans la demande
antérieure ni dans la demande internationale.\newline

\textbf{Conclusion: ne pas introduire ces dessins,
courriel d'avertissement pour modifications revendications
et description.}

\end{frame}

%--------------------------
\begin{frame}{Cas pratique}

\begin{block}{Conclusions}
\begin{itemize}
  \item ouvrir la phase régionale EP avant fin mai 2025:\\
  déposer la requête d’entrée en phase et paiement de la taxe nationale
  \item demander rapidement la signature de FR-O ;
  \item attendre l’avis de l’OEB sur l’unité, réponse selon option A ou B ;
  \item avertir sur l’introduction de matière nouvelle.
\end{itemize}
\end{block}

\end{frame}
